\subsection{サイクルがないこと}

\begin{theorem}
  \label{thm:no_cycle}
  \lean{no_cycle}
  \leanok
  \uses{}
  $C$はGCMであるとする．
  $k \in \N,\ 3 \le k$とし，$f : J := \{1, \ldots, k\} \hookrightarrow I$とする．
  任意の$i \in J$に対し，$C_{f(i), f(i+1 \bmod k)} \ne 0$とする．
  すなわち，頂点$f(1), \ldots, f(k)$がサイクルであるとする．
  このとき，$S(C)$は正定値でない．
\end{theorem}

\begin{proof}
  $x := \transpose{(x_1, x_2, \ldots, x_n)} \in \R^n$を以下で定める：
  \begin{equation}
    x_i = \begin{cases}
      1, & \exists j \in J \st i = f(j),\\
      0, & \textrm{otherwise}.
    \end{cases}
  \end{equation}
  $x_{f(0)} = 1$であるから$x \ne 0$であり，定義から明らかに$0 \le x_i$である．

  補題~\ref{lem:quadform_nonpos_of_neighbor_bound}が使いたい．
  そこで，$0 < x_i$として，以下を示す：
  \begin{equation}
    2 x_i \le \sum_{j \in N(i)} x_j.
  \end{equation}
  今，$x_i \ne 0$であるから，$x$の定義より$i = f(r)$を満たす$r \in J$が取れる．
  サイクル上にない頂点に関しては$x_* = 0$であるから，$i = f(r)$の近傍にあり，かつサイクル上にある頂点のみを考えれば良い．
  すなわち，
  \begin{equation}
    2 x_i
    = 2
    \le \#\set{j \in N(f(r))}{\exists u \st j = f(u)}
    = \sum_{j \in N(i)} x_j
  \end{equation}
  を示せばよい．
  \begin{equation}
    r_1 = r + 1 \bmod k,\quad r_2 = r + k - 1 \bmod k
  \end{equation}
  とする(サイクル上で$r$に隣接する2つの頂点である)．
  \begin{equation}
    r \ne r_1,\quad r \ne r_2,\quad r_1 \ne r_2
  \end{equation}
  である．
  \begin{equation}
    \{f(r_1), f(r_2)\} \subseteq \set{j \in N(f(r))}{\exists u \st j = f(u)}
  \end{equation}
  を示す．
  $f(r_i) \in N(f(r))$，すなわち$C_{f(r), f(r_i)} \ne 0$かつ$r \ne r_i$を示せば十分である．
  前者は定理の仮定と$r_i$の定め方からわかり，後者は先ほど示した通りである．
  また，単射性から$f(r_1) \ne f(r_2)$であるので，$\#\{f(r_1), f(r_2)\} = 2$もわかる．
  
  以上より，補題~\ref{lem:quadform_nonpos_of_neighbor_bound}の仮定をすべて満たすから
  \begin{equation}
    \transpose{x} S(C) x \le 0
  \end{equation}
  が成り立つ．
  ゆえに，$S(C)$は正定値でない．
\end{proof}